\documentclass[aspectratio=169,10pt]{beamer}

% Shared Beamer setup for Fall 2026 course slides.
% Clean Cal Poly-inspired theme: no custom class files or uncommon packages.
\mode<presentation>{
  \usetheme{default}
  \usefonttheme{professionalfonts}
}

\definecolor{CalPolyGreen}{HTML}{154734}
\definecolor{CalPolyGold}{HTML}{BD8B13}
\definecolor{CalPolySage}{HTML}{7FA28F}
\definecolor{CalPolyMint}{HTML}{DDF1E7}
\definecolor{CalPolySand}{HTML}{A29F86}
\definecolor{CalPolyBeige}{HTML}{EFEEDE}
\definecolor{DeckInk}{HTML}{1F2933}
\definecolor{DeckMuted}{HTML}{5F6B7A}
\definecolor{DeckSoft}{HTML}{F7F7F5}

\setbeamersize{text margin left=0.55in,text margin right=0.55in}
\setbeamercolor{normal text}{fg=DeckInk,bg=white}
\setbeamercolor{structure}{fg=CalPolyGreen}
\setbeamercolor{title}{fg=white}
\setbeamercolor{frametitle}{fg=CalPolyGreen,bg=white}
\setbeamercolor{block title}{bg=CalPolySage,fg=white}
\setbeamercolor{block body}{bg=CalPolyMint,fg=black}
\setbeamercolor{alerted text}{fg=CalPolyGold}

\setbeamerfont{title}{size=\Huge,series=\bfseries}
\setbeamerfont{frametitle}{size=\Large,series=\bfseries}
\setbeamerfont{block title}{size=\normalsize,series=\bfseries}

\setbeamertemplate{navigation symbols}{}
\setbeamertemplate{itemize item}{\color{CalPolyGold}$\blacktriangleright$}
\setbeamertemplate{itemize subitem}{\color{CalPolyGreen}$\bullet$}
\setbeamertemplate{footline}{
  \hfill{\color{DeckMuted}\scriptsize\insertframenumber/\inserttotalframenumber}
  \hspace{0.18in}\vspace{0.12in}
}
\setbeamertemplate{frametitle}{
  \vspace{0.18in}
  \hspace{0.02in}\insertframetitle\par
  \vspace{0.08in}
  \noindent\makebox[\linewidth][l]{%
    {\color{CalPolyGold}\rule{0.55in}{2pt}}%
    {\color{CalPolyGreen}\rule{\dimexpr\linewidth-0.55in\relax}{2pt}}%
  }
}

\newcommand{\CourseNumber}{Course}
\newcommand{\CourseTitle}{Course Title}
\newcommand{\LectureNumber}{0}
\newcommand{\LectureTitle}{Lecture Title}
\newcommand{\LectureDate}{Date}
\newcommand{\CourseUnit}{Unit}

\newcommand{\coursenumber}[1]{\renewcommand{\CourseNumber}{#1}}
\newcommand{\coursetitle}[1]{\renewcommand{\CourseTitle}{#1}}
\newcommand{\lecturenumber}[1]{\renewcommand{\LectureNumber}{#1}}
\newcommand{\lecturetitle}[1]{\renewcommand{\LectureTitle}{#1}}
\newcommand{\lecturedate}[1]{\renewcommand{\LectureDate}{#1}}
\newcommand{\courseunit}[1]{\renewcommand{\CourseUnit}{#1}}

\author{Kirk Duran}
\institute{California Polytechnic State University, San Luis Obispo}
\date{\LectureDate}

\newcommand{\makedecktitle}{
  \begingroup
  \setbeamercolor{background canvas}{bg=CalPolyGreen}
  \begin{frame}[plain]
    \vfill
    \begin{center}
      {\usebeamerfont{title}\color{white}\LectureTitle\par}
      \vspace{0.18in}
      {\Large\color{white}Lecture \LectureNumber\par}
    \end{center}
    \vfill
    \noindent{\color{white}\rule{\linewidth}{1.2pt}}\par
    \vspace{0.08in}
    \noindent\colorbox{CalPolyGold}{\makebox[0.24\linewidth][c]{\color{white}\Large\CourseNumber}}
    \hspace{0.12in}{\color{white}\Large Kirk Duran}
  \end{frame}
  \endgroup
}

\newenvironment{keyidea}{\begin{block}{Key idea}}{\end{block}}
\newenvironment{activity}{\begin{block}{In class}}{\end{block}}
\newenvironment{cpdefinition}{\begin{block}{Definition}}{\end{block}}
\newenvironment{exampleblockcp}{
  \begingroup
  \setbeamercolor{block title}{bg=CalPolySand,fg=white}
  \setbeamercolor{block body}{bg=CalPolyBeige,fg=black}
  \begin{block}{Example}
}{
  \end{block}
  \endgroup
}

\newcommand{\imageplaceholder}[2][0.3in]{%
  \noindent\makebox[\linewidth][c]{%
    \fcolorbox{CalPolySage}{CalPolyBeige}{%
      \parbox[c][#1][c]{0.86\linewidth}{%
        \centering
        {\color{DeckMuted}\scriptsize Image placeholder: }%
        {\color{DeckInk}\scriptsize\ttfamily\detokenize{#2}}%
      }%
    }%
  }%
}

\newcommand{\sectionframe}[1]{
  \begingroup
  \setbeamercolor{background canvas}{bg=CalPolyGreen}
  \begin{frame}[plain]
    \vfill
    {\color{white}\LARGE\bfseries #1\par}
    \vspace{0.18in}
    {\color{CalPolyGold}\rule{0.22\paperwidth}{3pt}}
    {\color{white}\rule{0.62\paperwidth}{3pt}}
    \vfill
  \end{frame}
  \endgroup
}

\newcommand{\placeholdernote}{
  \begin{block}{Placeholder}
    Replace these bullets with the final lecture material, examples, and in-class activities.
  \end{block}
}


\pdfstringdefDisableCommands{\def\translate#1{#1}}

\graphicspath{{../assets/lecture-18/}}

% If an image file is not yet available, the slide displays a labeled
% placeholder using the intended filename. Place the matching file in
% ../assets/lecture-18/ to replace the placeholder automatically.
\newcommand{\lectureimage}[2][1.75in]{%
  \IfFileExists{../assets/lecture-18/#2}{%
    \includegraphics[width=\linewidth,height=#1,keepaspectratio]{#2}%
  }{%
    \fbox{%
      \parbox[c][#1][c]{0.94\linewidth}{%
        \centering
        \small\textbf{Image Placeholder}\\[0.08in]
        \scriptsize\texttt{\detokenize{#2}}
      }%
    }%
  }%
}

\setbeamersize{text margin left=0.42in,text margin right=0.42in}

\coursenumber{CS 1001}
\coursetitle{Introduction to Computing}
\lecturenumber{18}
\lecturetitle{Lists II --- Slicing and Built-ins}
\lecturedate{Fall 2026}
\courseunit{7. Sequential Data}

\begin{document}

\makedecktitle

\begin{frame}{Today}
  \begin{itemize}
    \item Extract portions of a list using slicing.
    \item Interpret \texttt{start}, \texttt{stop}, and the optional \texttt{step}.
    \item Distinguish operations that create a new list from operations that mutate an existing list.
    \item Measure list length with \texttt{len()}.
    \item Add elements using \texttt{append()}, \texttt{extend()}, concatenation, and \texttt{insert()}.
    \item Remove elements using \texttt{pop()} and \texttt{remove()}.
    \item Refine the list memory model to account for changing list size.
  \end{itemize}

  \begin{keyidea}
    A Python list behaves like a resizable sequence: positions remain ordered while Python manages the storage required when the sequence grows or shrinks.
  \end{keyidea}
\end{frame}

\sectionframe{Part 1: Slicing a Sequence}

\begin{frame}{Indexing Selects One Element; Slicing Selects a Range}
  \begin{columns}[T,onlytextwidth]
    \begin{column}{0.47\textwidth}
      \begin{block}{Indexing}
        \centering
        \texttt{values[2]} $\longrightarrow$ \texttt{30}
      \end{block}
      \begin{itemize}
        \item Selects one position.
        \item Evaluates to one element value.
      \end{itemize}
    \end{column}
    \begin{column}{0.47\textwidth}
      \begin{block}{Slicing}
        \centering
        \texttt{values[1:4]} $\longrightarrow$ \texttt{[20, 30, 40]}
      \end{block}
      \begin{itemize}
        \item Selects a sequence of positions.
        \item Evaluates to a new list.
      \end{itemize}
    \end{column}
  \end{columns}

  \begin{keyidea}
    Indexing asks for one element. Slicing asks for a subsequence.
  \end{keyidea}
\end{frame}

\begin{frame}{Definition: Slice}
  \begin{cpdefinition}
    A \textbf{slice} specifies a range of positions from a sequence and produces a new sequence containing the selected elements.
  \end{cpdefinition}

  \begin{block}{General list-slice syntax}
    \centering
    \Large\texttt{list[start:stop:step]}
  \end{block}

  \begin{itemize}
    \item \texttt{start}: first selected index.
    \item \texttt{stop}: boundary index at which selection stops.
    \item \texttt{step}: distance between selected indexes.
    \item Any of these components may be omitted when its default behavior is desired.
  \end{itemize}
\end{frame}

\begin{frame}{The Stop Index Is Exclusive}
  \begin{block}{Example}
    \texttt{values = [10, 20, 30, 40, 50]}\\[0.10in]
    \centering
    \Large\texttt{values[1:4]} $\longrightarrow$ \texttt{[20, 30, 40]}
  \end{block}

  \begin{center}
    \renewcommand{\arraystretch}{1.25}
    \begin{tabular}{c|c|c|c|c|c}
      index & 0 & \textbf{1} & \textbf{2} & \textbf{3} & 4 \\\hline
      value & 10 & \textbf{20} & \textbf{30} & \textbf{40} & 50
    \end{tabular}
  \end{center}

  \begin{keyidea}
    \texttt{start} is included. \texttt{stop} is not included. The slice \texttt{1:4} selects indexes \texttt{1}, \texttt{2}, and \texttt{3}.
  \end{keyidea}
\end{frame}





\begin{frame}{Omitted Slice Bounds Use Sequence Boundaries}
  \begin{block}{Examples}
    \texttt{values = [10, 20, 30, 40, 50]}\\[0.08in]
    \texttt{values[:3]} $\longrightarrow$ \texttt{[10, 20, 30]}\\
    \texttt{values[2:]} $\longrightarrow$ \texttt{[30, 40, 50]}\\
    \texttt{values[:]} $\longrightarrow$ \texttt{[10, 20, 30, 40, 50]}
  \end{block}

  \begin{itemize}
    \item With a positive step, omitted \texttt{start} begins at the first element.
    \item Omitted \texttt{stop} continues through the final element.
    \item Omitting both bounds selects the complete list and produces a new list.
  \end{itemize}
\end{frame}

\begin{frame}{Slice Bounds May Extend Beyond the Valid Index Range}
  \begin{block}{Examples}
    \texttt{values = [10, 20, 30, 40]}\\[0.08in]
    \texttt{values[1:100]} $\longrightarrow$ \texttt{[20, 30, 40]}\\
    \texttt{values[100:200]} $\longrightarrow$ \texttt{[]}
  \end{block}

  \begin{itemize}
    \item Ordinary indexing such as \texttt{values[100]} raises \texttt{IndexError}.
    \item Slicing adjusts an out-of-range boundary to the available sequence.
    \item A valid slice may therefore produce an empty list.
  \end{itemize}
\end{frame}

\begin{frame}{Negative Indexes Also Work in Slice Bounds}
  \begin{block}{Example}
    \texttt{values = [10, 20, 30, 40, 50]}\\[0.10in]
    \texttt{values[1:-1]} $\longrightarrow$ \texttt{[20, 30, 40]}
  \end{block}

  \begin{itemize}
    \item \texttt{-1} identifies the final element's position.
    \item Because \texttt{stop} is exclusive, the element at index \texttt{-1} is not selected.
    \item The same index rules from Lecture 17 apply inside slice notation.
  \end{itemize}
\end{frame}

\sectionframe{Part 2: The Optional Step}

\begin{frame}{The Step Controls the Distance Between Selected Indexes}
  \begin{block}{Example}
    \texttt{values = [10, 20, 30, 40, 50, 60]}\\[0.10in]
    \texttt{values[0:6:2]} $\longrightarrow$ \texttt{[10, 30, 50]}
  \end{block}

  \begin{center}
    \renewcommand{\arraystretch}{1.2}
    \begin{tabular}{c|c|c|c}
      selected index & 0 & 2 & 4 \\\hline
      selected value & 10 & 30 & 50
    \end{tabular}
  \end{center}

  \begin{itemize}
    \item If \texttt{step} is omitted, it defaults to \texttt{1}.
    \item A step of \texttt{0} is invalid because the index would never advance; Python raises \texttt{ValueError}.
  \end{itemize}
\end{frame}





\begin{frame}{A Negative Step Selects Positions in Decreasing Index Order}
  \begin{block}{Example}
    \texttt{values = [10, 20, 30, 40, 50]}\\[0.10in]
    \texttt{values[4:1:-1]} $\longrightarrow$ \texttt{[50, 40, 30]}
  \end{block}

  \begin{itemize}
    \item Start at index \texttt{4}.
    \item Subtract \texttt{1} after each selected position.
    \item Stop before index \texttt{1}.
  \end{itemize}

  \begin{keyidea}
    The stop boundary remains exclusive even when the direction is reversed.
  \end{keyidea}
\end{frame}

\begin{frame}{The Slice \texttt{[::-1]} Produces a Reversed Copy}
  \begin{columns}[T,onlytextwidth]
    \begin{column}{0.50\textwidth}
      \begin{block}{Code}
        \texttt{values = [10, 20, 30, 40]}\\[0.08in]
        \texttt{reversed\_values = values[::-1]}
      \end{block}

      \begin{block}{Result}
        \texttt{reversed\_values}\
        $\longrightarrow$ \texttt{[40, 30, 20, 10]}
      \end{block}
    \end{column}
    \begin{column}{0.44\textwidth}
      % Intended visual: positive indexes left-to-right and arrows moving right-to-left.
      \lectureimage[2.25in]{negative-step-reverse-slice.png}
    \end{column}
  \end{columns}

  \begin{alertblock}{Important}
    The original list is not mutated by the slice.
  \end{alertblock}
\end{frame}

\begin{frame}{Practice: Predict the Slice Results}
  \begin{block}{Given}
    \centering
    \texttt{data = [5, 10, 15, 20, 25, 30]}
  \end{block}

  \begin{activity}
    Predict each result before evaluating it in Python.
    \begin{enumerate}
      \item \texttt{data[1:4]}
      \item \texttt{data[:3]}
      \item \texttt{data[3:]}
      \item \texttt{data[::2]}
      \item \texttt{data[4:1:-1]}
      \item \texttt{data[::-1]}
    \end{enumerate}
  \end{activity}
\end{frame}

\sectionframe{Part 3: Length and Operation Categories}

\begin{frame}{\texttt{len()} Returns the Number of Elements}
  \begin{cpdefinition}
    The built-in function \texttt{len(sequence)} returns the number of elements in the sequence as an integer.
  \end{cpdefinition}

  \begin{block}{Examples}
    \texttt{len([10, 20, 30, 40])} $\longrightarrow$ \texttt{4}\\
    \texttt{len([])} $\longrightarrow$ \texttt{0}
  \end{block}

  \begin{keyidea}
    Length counts elements. It is not an index.
  \end{keyidea}
\end{frame}

\begin{frame}{Length and the Last Positive Index Differ by One}
  \begin{block}{Example}
    \texttt{values = [10, 20, 30, 40]}\\[0.08in]
    \texttt{len(values)} $\longrightarrow$ \texttt{4}\\
    \texttt{values[3]} $\longrightarrow$ \texttt{40}
  \end{block}

  \begin{itemize}
    \item Four elements occupy positive indexes \texttt{0}, \texttt{1}, \texttt{2}, and \texttt{3}.
    \item Therefore, the final valid positive index is \texttt{len(values) - 1}.
    \item \texttt{values[len(values)]} would be one position beyond the list and raises \texttt{IndexError}.
  \end{itemize}
\end{frame}

\begin{frame}{Built-in Functions and List Methods Use Different Syntax}
  \begin{columns}[T,onlytextwidth]
    \begin{column}{0.47\textwidth}
      \begin{block}{Built-in function}
        \centering
        \Large\texttt{len(values)}
      \end{block}
      \begin{itemize}
        \item The list is passed as an argument.
        \item The call returns a result.
      \end{itemize}
    \end{column}
    \begin{column}{0.47\textwidth}
      \begin{block}{List method}
        \centering
        \Large\texttt{values.append(40)}
      \end{block}
      \begin{itemize}
        \item The method is invoked through the list value.
        \item Many list methods mutate that list.
      \end{itemize}
    \end{column}
  \end{columns}
\end{frame}



\sectionframe{Part 4: Growing a List}

\begin{frame}{\texttt{append()} Adds One Element at the End}
  \begin{block}{Example}
    \texttt{numbers = [10, 20, 30]}\\
    \texttt{numbers.append(40)}\\[0.08in]
    \texttt{numbers} $\longrightarrow$ \texttt{[10, 20, 30, 40]}
  \end{block}

  \begin{itemize}
    \item The existing list is mutated.
    \item Its logical length increases by one.
    \item The supplied argument becomes exactly one new element.
    \item \texttt{append()} returns \texttt{None}.
  \end{itemize}
\end{frame}

\begin{frame}{Growing a List May Require Resizing Its Storage}
  \begin{columns}[T,onlytextwidth]
    \begin{column}{0.54\textwidth}
      \begin{itemize}
        \item In our working model, adding an element requires another adjacent slot.
        \item Python manages this storage change automatically.
        \item In CPython, the list uses a contiguous array of object references with some internal capacity.
        \item If growth exceeds available capacity, Python may allocate a larger reference array and copy the existing references.
      \end{itemize}
    \end{column}
    \begin{column}{0.40\textwidth}
      \lectureimage[2.30in]{append-resize-working-model.png}
    \end{column}
  \end{columns}

  \begin{keyidea}
    Dynamic resizing is hidden behind the list abstraction; the programmer specifies the sequence operation, not the memory-management procedure.
  \end{keyidea}
\end{frame}

\begin{frame}{\texttt{extend()} Adds the Elements of Another List}
  \begin{block}{Example}
    \texttt{numbers = [10, 20, 30]}\\
    \texttt{numbers.extend([40, 50])}\\[0.08in]
    \texttt{numbers} $\longrightarrow$ \texttt{[10, 20, 30, 40, 50]}
  \end{block}

  \begin{itemize}
    \item The original list is mutated.
    \item Each element from the supplied list becomes a separate element at the end.
    \item The logical length increases by the number of added elements.
    \item \texttt{extend()} returns \texttt{None}.
  \end{itemize}
\end{frame}

\begin{frame}{\texttt{append()} and \texttt{extend()} Are Not Equivalent}
  \begin{columns}[T,onlytextwidth]
    \begin{column}{0.47\textwidth}
      \begin{block}{Append one value}
        \texttt{a = [10, 20]}\\
        \texttt{a.append([30, 40])}\\[0.08in]
        \texttt{a} $\longrightarrow$\\
        \texttt{[10, 20, [30, 40]]}
      \end{block}
    \end{column}
    \begin{column}{0.47\textwidth}
      \begin{block}{Extend by values}
        \texttt{b = [10, 20]}\\
        \texttt{b.extend([30, 40])}\\[0.08in]
        \texttt{b} $\longrightarrow$\\
        \texttt{[10, 20, 30, 40]}
      \end{block}
    \end{column}
  \end{columns}

  \begin{keyidea}
    \texttt{append(x)} adds \texttt{x} as one element. \texttt{extend(other)} adds the elements contained in \texttt{other}.
  \end{keyidea}
\end{frame}

\begin{frame}{Concatenation Combines Lists Without Mutating Either Operand}
  \begin{block}{Example}
    \texttt{left = [10, 20]}\\
    \texttt{right = [30, 40]}\\
    \texttt{combined = left + right}
  \end{block}

  \begin{block}{State afterward}
    \texttt{left} $\longrightarrow$ \texttt{[10, 20]}\\
    \texttt{right} $\longrightarrow$ \texttt{[30, 40]}\\
    \texttt{combined} $\longrightarrow$ \texttt{[10, 20, 30, 40]}
  \end{block}

  \begin{keyidea}
    Concatenation creates a new list; \texttt{extend()} changes an existing list. Similar visible contents can arise from different state-change semantics.
  \end{keyidea}
\end{frame}



\begin{frame}{\texttt{insert()} Adds One Element at a Specific Position}
  \begin{block}{Example}
    \texttt{numbers = [10, 20, 40, 50]}\\
    \texttt{numbers.insert(2, 30)}\\[0.08in]
    \texttt{numbers} $\longrightarrow$ \texttt{[10, 20, 30, 40, 50]}
  \end{block}

  \begin{itemize}
    \item The first argument specifies the insertion position.
    \item The second argument is the value to insert.
    \item The existing element at that position and later elements move right.
    \item \texttt{insert()} mutates the list and returns \texttt{None}.
  \end{itemize}
\end{frame}



\sectionframe{Part 5: Removing Elements}

\begin{frame}{\texttt{pop()} Removes by Position and Returns the Removed Element}
  \begin{block}{Example}
    \texttt{numbers = [10, 20, 30, 40]}\\
    \texttt{removed = numbers.pop(1)}
  \end{block}

  \begin{block}{State afterward}
    \texttt{removed} $\longrightarrow$ \texttt{20}\\
    \texttt{numbers} $\longrightarrow$ \texttt{[10, 30, 40]}
  \end{block}

  \begin{itemize}
    \item With no index argument, \texttt{pop()} removes and returns the final element.
    \item An out-of-range supplied index raises \texttt{IndexError}.
  \end{itemize}
  \begin{keyidea}
    \texttt{pop()} both mutates the list and produces the removed element as its return value.
  \end{keyidea}
\end{frame}



\begin{frame}{Removing from the Middle Shifts Later Positions Left}
  \begin{center}
    % Intended visual: [10][20][30][40] with 20 removed and 30/40 shifting left.
    \lectureimage[2.10in]{pop-shift-left.png}
  \end{center}

  \begin{itemize}
    \item Removing a middle element creates a gap in the ordered positions.
    \item Later elements shift left so the list continues to use consecutive indexes starting at \texttt{0}.
    \item The logical length decreases by one.
  \end{itemize}
\end{frame}

\begin{frame}{\texttt{remove()} Removes by Value}
  \begin{block}{Example}
    \texttt{fruits = ["apple", "banana", "cherry", "banana"]}\\
    \texttt{fruits.remove("banana")}\\[0.08in]
    \texttt{fruits} $\longrightarrow$ \texttt{["apple", "cherry", "banana"]}
  \end{block}

  \begin{itemize}
    \item \texttt{remove(value)} searches for the first equal value.
    \item That first matching element is removed.
    \item The method mutates the list and returns \texttt{None}.
    \item If no matching value exists, Python raises \texttt{ValueError}.
  \end{itemize}
\end{frame}



\sectionframe{Part 6: Reversal, Resizing, and Tracing}



\begin{frame}{Reversed Slice Versus \texttt{reverse()}}
  \begin{columns}[T,onlytextwidth]
    \begin{column}{0.47\textwidth}
      \begin{block}{New reversed list}
        \texttt{a = [1, 2, 3]}\\
        \texttt{b = a[::-1]}\\[0.08in]
        \texttt{a} $\rightarrow$ \texttt{[1, 2, 3]}\\
        \texttt{b} $\rightarrow$ \texttt{[3, 2, 1]}
      \end{block}
    \end{column}
    \begin{column}{0.47\textwidth}
      \begin{block}{Mutate existing list}
        \texttt{a = [1, 2, 3]}\\
        \texttt{a.reverse()}\\[0.08in]
        \texttt{a} $\rightarrow$ \texttt{[3, 2, 1]}
      \end{block}
    \end{column}
  \end{columns}

  \begin{keyidea}
    Similar visible results can have different mutation semantics.
  \end{keyidea}
\end{frame}

\begin{frame}{List Size Is Logical; Capacity Is an Implementation Detail}
  \begin{columns}[T,onlytextwidth]
    \begin{column}{0.54\textwidth}
      \begin{itemize}
        \item \texttt{len(values)} reports the logical number of elements.
        \item The underlying storage may reserve additional capacity for future growth.
        \item Adding or removing elements changes the logical sequence immediately.
        \item Any necessary allocation, copying, or capacity management is performed by Python.
      \end{itemize}
    \end{column}
    \begin{column}{0.40\textwidth}
      % Intended visual: logical length 4 inside a backing capacity of 6 or 8 reference slots.
      \lectureimage[2.35in]{logical-length-vs-capacity.png}
    \end{column}
  \end{columns}
\end{frame}



\begin{frame}{Trace State After Each Mutating Operation}
  \begin{block}{Code}
    \texttt{values = [10, 20, 30]}\\
    \texttt{values.append(40)}\\
    \texttt{values.insert(1, 15)}\\
    \texttt{removed = values.pop(3)}
  \end{block}

  \begin{center}
    \renewcommand{\arraystretch}{1.28}
    \begin{tabular}{l|l|l}
      \textbf{After statement} & \textbf{values} & \textbf{removed} \\\hline
      creation & \texttt{[10, 20, 30]} & --- \\
      \texttt{append(40)} & \texttt{[10, 20, 30, 40]} & --- \\
      \texttt{insert(1, 15)} & \texttt{[10, 15, 20, 30, 40]} & --- \\
      \texttt{pop(3)} & \texttt{[10, 15, 20, 40]} & \texttt{30}
    \end{tabular}
  \end{center}
\end{frame}

\begin{frame}{Trace New-List Operations Separately from Mutation}
  \begin{block}{Code}
    \texttt{values = [10, 20, 30, 40]}\\
    \texttt{middle = values[1:3]}\\
    \texttt{combined = values + [50, 60]}
  \end{block}

  \begin{center}
    \renewcommand{\arraystretch}{1.3}
    \begin{tabular}{l|l}
      \textbf{Name} & \textbf{Value afterward} \\\hline
      \texttt{values} & \texttt{[10, 20, 30, 40]} \\
      \texttt{middle} & \texttt{[20, 30]} \\
      \texttt{combined} & \texttt{[10, 20, 30, 40, 50, 60]}
    \end{tabular}
  \end{center}

  \begin{keyidea}
    Neither slicing nor concatenation changes \texttt{values} in this trace.
  \end{keyidea}
\end{frame}

\begin{frame}[shrink=4]{Operation Summary}
  \begin{center}
    \renewcommand{\arraystretch}{1.25}
    \begin{tabular}{l|c|l}
      \textbf{Operation} & \textbf{Mutates list?} & \textbf{Result of expression/call} \\\hline
      \texttt{a[start:stop:step]} & no & new list \\
      \texttt{len(a)} & no & integer length \\
      \texttt{a.append(x)} & yes & \texttt{None} \\
      \texttt{a.extend(b)} & yes & \texttt{None} \\
      \texttt{a + b} & no & new list \\
      \texttt{a.insert(i, x)} & yes & \texttt{None} \\
      \texttt{a.pop(i)} & yes & removed element \\
      \texttt{a.remove(x)} & yes & \texttt{None} \\
      \texttt{a.reverse()} & yes & \texttt{None}
    \end{tabular}
  \end{center}
\end{frame}

\begin{frame}{Integrated Prediction Exercise}
  \begin{block}{Code}
    \texttt{data = [5, 10, 15, 20]}\\
    \texttt{part = data[1:3]}\\
    \texttt{data.append(25)}\\
    \texttt{data.insert(2, 12)}\\
    \texttt{removed = data.pop()}\\
    \texttt{data.reverse()}
  \end{block}

  \begin{activity}
    Without running the program, determine the final values of \texttt{data}, \texttt{part}, and \texttt{removed}. For every statement, classify it as mutating or non-mutating.
  \end{activity}
\end{frame}

\begin{frame}{Takeaways}
  \begin{itemize}
    \item Slicing uses \texttt{start:stop:step}; the stop boundary is exclusive.
    \item Omitting slice components invokes defaults, and a negative step reverses the direction of selection.
    \item \texttt{len()} reports the number of elements, not the last index.
    \item \texttt{append()}, \texttt{extend()}, \texttt{insert()}, \texttt{pop()}, \texttt{remove()}, and \texttt{reverse()} mutate a list.
    \item Slicing and concatenation create new lists instead of mutating their operands.
    \item Python automatically manages storage as a list grows or shrinks; our programming model remains an ordered sequence of adjacent positions.
  \end{itemize}

  \begin{keyidea}
    When tracing list code, always track both the sequence contents and whether an operation mutates an existing list or produces a new value.
  \end{keyidea}
\end{frame}

\end{document}
